
\documentclass{article}
\usepackage{colortbl}
\usepackage{makecell}
\usepackage{multirow}
\usepackage{supertabular}

\begin{document}

\newcounter{utterance}

\centering \large Interaction Transcript for game `matchit\_ascii', experiment `similar\_grid\_2', episode 3 with Apertus{-}v1.5{-}8B.
\vspace{24pt}

{ \footnotesize  \setcounter{utterance}{1}
\setlength{\tabcolsep}{0pt}
\begin{supertabular}{c@{$\;$}|p{.15\linewidth}@{}p{.15\linewidth}p{.15\linewidth}p{.15\linewidth}p{.15\linewidth}p{.15\linewidth}}
    \# & $\;$A & \multicolumn{4}{c}{Game Master} & $\;\:$B\\
    \hline

    \theutterance \stepcounter{utterance}  
    & & \multicolumn{4}{p{0.6\linewidth}}{
        \cellcolor[rgb]{0.9,0.9,0.9}{
            \makecell[{{p{\linewidth}}}]{
                \texttt{\tiny{[P1$\langle$GM]}}
                \texttt{U neemt deel aan een gezamenlijk raadspel.} \\
\texttt{Het doel is om erachter te komen of dit raster en een andere afbeelding die alleen ik kan zien, hetzelfde zijn.} \\
\texttt{Beschrijf eerst uw raster. Daarna geef ik mijn beschrijving en kunnen we elkaar vragen stellen over de rasters om erachter te komen of ze hetzelfde zijn. Begin nu uw korte rasterbeschrijving met "BESCHRIJVING:" gevolgd door de beschrijving. Voeg niets anders toe. Dit is uw raster:} \\
\texttt{X X X X X} \\
\texttt{▢ ▢ ▢ ▢ X} \\
\texttt{▢ ▢ X X X} \\
\texttt{▢ ▢ ▢ ▢ X} \\
\texttt{▢ ▢ ▢ ▢ X} \\
            }
        }
    }
    & \\ \\

    \theutterance \stepcounter{utterance}  
    & \multicolumn{4}{p{0.6\linewidth}}{
        \cellcolor[rgb]{0.9,0.9,0.9}{
            \makecell[{{p{\linewidth}}}]{
                \texttt{\tiny{[P1$\rangle$GM]}}
                \texttt{BESCHRIJVING: Een raster van vier rijen en vijf kolommen, waarbij de eerste rij bestaat uit vijf X{-}tekens. De tweede en vierde rij bevatten drie ▢{-}symbolen aan het begin, gevolgd door drie X{-}tekens in het midden en eindigen met één ▢. De derde rij begint met twee ▢{-}symbolen, gevolgd door drie X{-}tekens en eindigt weer met één ▢.} \\
            }
        }
    }
    & & \\ \\

    \theutterance \stepcounter{utterance}  
    & & & \multicolumn{2}{p{0.3\linewidth}}{
        \cellcolor[rgb]{0.9,0.9,0.9}{
            \makecell[{{p{\linewidth}}}]{
                \texttt{\tiny{[GM$|$GM]}}
                \texttt{continue} \\
            }
        }
    }
    & & \\ \\

    \theutterance \stepcounter{utterance}  
    & & \multicolumn{4}{p{0.6\linewidth}}{
        \cellcolor[rgb]{0.9,0.9,0.9}{
            \makecell[{{p{\linewidth}}}]{
                \texttt{\tiny{[GM$\rangle$P2]}}
                \texttt{U neemt deel aan een gezamenlijk raadspel.} \\
\texttt{Het doel is om erachter te komen of dit raster en een andere afbeelding die alleen ik kan zien, hetzelfde zijn.} \\
\texttt{Beschrijf eerst uw raster. Daarna geef ik mijn beschrijving en kunnen we elkaar vragen stellen over de rasters om erachter te komen of ze hetzelfde zijn. Begin nu uw korte rasterbeschrijving met "BESCHRIJVING:" gevolgd door de beschrijving. Voeg niets anders toe. Dit is uw raster:} \\
\texttt{X X X X X} \\
\texttt{▢ ▢ ▢ ▢ X} \\
\texttt{▢ X X X X} \\
\texttt{▢ ▢ ▢ ▢ X} \\
\texttt{▢ ▢ X ▢ X} \\
            }
        }
    }
    & \\ \\

    \theutterance \stepcounter{utterance}  
    & & & \multicolumn{4}{p{0.6\linewidth}}{
        \cellcolor[rgb]{0.9,0.9,0.9}{
            \makecell[{{p{\linewidth}}}]{
                \texttt{\tiny{[GM$\langle$P2]}}
                \texttt{BESCHRIJVING: Het raster bestaat uit vijf rijen en vijf kolommen. De eerste rij bevat vijf X{-}tekens. De tweede rij begint met drie lege vakken, gevolgd door vier X{-}tekens. De derde rij begint met één leeg vak en vervolgens drie X{-}tekens. De vierde rij lijkt op de tweede rij, met drie lege vakken en vier X{-}tekens. De vijfde rij begint met twee lege vakken, gevolgd door een X, een leeg vak, een X en eindigt met een X.} \\
            }
        }
    }
     \\ \\

    \theutterance \stepcounter{utterance}  
    & & & \multicolumn{2}{p{0.3\linewidth}}{
        \cellcolor[rgb]{0.9,0.9,0.9}{
            \makecell[{{p{\linewidth}}}]{
                \texttt{\tiny{[GM$|$GM]}}
                \texttt{continue} \\
            }
        }
    }
    & & \\ \\

    \theutterance \stepcounter{utterance}  
    & & \multicolumn{4}{p{0.6\linewidth}}{
        \cellcolor[rgb]{0.9,0.9,0.9}{
            \makecell[{{p{\linewidth}}}]{
                \texttt{\tiny{[GM$\rangle$P2]}}
                \texttt{Dit is mijn BESCHRIJVINGBESCHRIJVING: Een raster van vier rijen en vijf kolommen, waarbij de eerste rij bestaat uit vijf X{-}tekens. De tweede en vierde rij bevatten drie ▢{-}symbolen aan het begin, gevolgd door drie X{-}tekens in het midden en eindigen met één ▢. De derde rij begint met twee ▢{-}symbolen, gevolgd door drie X{-}tekens en eindigt weer met één ▢.} \\
\texttt{Stel nu een vraag om nieuwe aspecten van mijn raster te ontdekken die anders kunnen zijn dan jouw raster. Begin met "VRAAG:" en voeg niets anders toe.} \\
            }
        }
    }
    & \\ \\

    \theutterance \stepcounter{utterance}  
    & & & \multicolumn{4}{p{0.6\linewidth}}{
        \cellcolor[rgb]{0.9,0.9,0.9}{
            \makecell[{{p{\linewidth}}}]{
                \texttt{\tiny{[GM$\langle$P2]}}
                \texttt{VRAAG: Zijn er verschillen in het aantal lege vakken (▢) in de rijen waar je drie X{-}tekens noemt?} \\
            }
        }
    }
     \\ \\

    \theutterance \stepcounter{utterance}  
    & & & \multicolumn{2}{p{0.3\linewidth}}{
        \cellcolor[rgb]{0.9,0.9,0.9}{
            \makecell[{{p{\linewidth}}}]{
                \texttt{\tiny{[GM$|$GM]}}
                \texttt{continue} \\
            }
        }
    }
    & & \\ \\

    \theutterance \stepcounter{utterance}  
    & & \multicolumn{4}{p{0.6\linewidth}}{
        \cellcolor[rgb]{0.9,0.9,0.9}{
            \makecell[{{p{\linewidth}}}]{
                \texttt{\tiny{[P1$\langle$GM]}}
                \texttt{Dit is mijn BESCHRIJVINGBESCHRIJVING: Het raster bestaat uit vijf rijen en vijf kolommen. De eerste rij bevat vijf X{-}tekens. De tweede rij begint met drie lege vakken, gevolgd door vier X{-}tekens. De derde rij begint met één leeg vak en vervolgens drie X{-}tekens. De vierde rij lijkt op de tweede rij, met drie lege vakken en vier X{-}tekens. De vijfde rij begint met twee lege vakken, gevolgd door een X, een leeg vak, een X en eindigt met een X.} \\
\texttt{VRAAG: Zijn er verschillen in het aantal lege vakken (▢) in de rijen waar je drie X{-}tekens noemt?Begin je antwoord met "ANTWOORD:" en voeg niets anders toe.} \\
            }
        }
    }
    & \\ \\

    \theutterance \stepcounter{utterance}  
    & \multicolumn{4}{p{0.6\linewidth}}{
        \cellcolor[rgb]{0.9,0.9,0.9}{
            \makecell[{{p{\linewidth}}}]{
                \texttt{\tiny{[P1$\rangle$GM]}}
                \texttt{ANTWOORD: Nee, in de rijen waar drie X{-}tekens worden genoemd, zijn er allemaal drie lege vakken (▢).} \\
            }
        }
    }
    & & \\ \\

    \theutterance \stepcounter{utterance}  
    & & & \multicolumn{2}{p{0.3\linewidth}}{
        \cellcolor[rgb]{0.9,0.9,0.9}{
            \makecell[{{p{\linewidth}}}]{
                \texttt{\tiny{[GM$|$GM]}}
                \texttt{continue} \\
            }
        }
    }
    & & \\ \\

    \theutterance \stepcounter{utterance}  
    & & \multicolumn{4}{p{0.6\linewidth}}{
        \cellcolor[rgb]{0.9,0.9,0.9}{
            \makecell[{{p{\linewidth}}}]{
                \texttt{\tiny{[P1$\langle$GM]}}
                \texttt{Stel nu een vraag om nieuwe aspecten van mijn raster te ontdekken die anders kunnen zijn dan jouw raster. Begin met "VRAAG:" en voeg niets anders toe.} \\
            }
        }
    }
    & \\ \\

    \theutterance \stepcounter{utterance}  
    & \multicolumn{4}{p{0.6\linewidth}}{
        \cellcolor[rgb]{0.9,0.9,0.9}{
            \makecell[{{p{\linewidth}}}]{
                \texttt{\tiny{[P1$\rangle$GM]}}
                \texttt{VRAAG: Zijn er X{-}tekens in kolommen die niet overeenkomen met de X{-}tekens in jouw raster?} \\
            }
        }
    }
    & & \\ \\

    \theutterance \stepcounter{utterance}  
    & & & \multicolumn{2}{p{0.3\linewidth}}{
        \cellcolor[rgb]{0.9,0.9,0.9}{
            \makecell[{{p{\linewidth}}}]{
                \texttt{\tiny{[GM$|$GM]}}
                \texttt{continue} \\
            }
        }
    }
    & & \\ \\

    \theutterance \stepcounter{utterance}  
    & & \multicolumn{4}{p{0.6\linewidth}}{
        \cellcolor[rgb]{0.9,0.9,0.9}{
            \makecell[{{p{\linewidth}}}]{
                \texttt{\tiny{[GM$\rangle$P2]}}
                \texttt{ANTWOORD: Nee, in de rijen waar drie X{-}tekens worden genoemd, zijn er allemaal drie lege vakken (▢).} \\
\texttt{VRAAG: Zijn er X{-}tekens in kolommen die niet overeenkomen met de X{-}tekens in jouw raster?Begin je antwoord met "ANTWOORD:" en voeg niets anders toe.} \\
            }
        }
    }
    & \\ \\

    \theutterance \stepcounter{utterance}  
    & & & \multicolumn{4}{p{0.6\linewidth}}{
        \cellcolor[rgb]{0.9,0.9,0.9}{
            \makecell[{{p{\linewidth}}}]{
                \texttt{\tiny{[GM$\langle$P2]}}
                \texttt{ANTWOORD: Ik kan geen kolommen met X{-}tekens visualiseren, aangezien ik alleen de rasterbeschrijving heb. Ik kan dus geen antwoord geven op vragen over kolommen zonder verdere informatie.} \\
            }
        }
    }
     \\ \\

    \theutterance \stepcounter{utterance}  
    & & & \multicolumn{2}{p{0.3\linewidth}}{
        \cellcolor[rgb]{0.9,0.9,0.9}{
            \makecell[{{p{\linewidth}}}]{
                \texttt{\tiny{[GM$|$GM]}}
                \texttt{continue} \\
            }
        }
    }
    & & \\ \\

    \theutterance \stepcounter{utterance}  
    & & \multicolumn{4}{p{0.6\linewidth}}{
        \cellcolor[rgb]{0.9,0.9,0.9}{
            \makecell[{{p{\linewidth}}}]{
                \texttt{\tiny{[GM$\rangle$P2]}}
                \texttt{Stel nu een vraag om nieuwe aspecten van mijn raster te ontdekken die anders kunnen zijn dan jouw raster. Begin met "VRAAG:" en voeg niets anders toe.} \\
            }
        }
    }
    & \\ \\

    \theutterance \stepcounter{utterance}  
    & & & \multicolumn{4}{p{0.6\linewidth}}{
        \cellcolor[rgb]{0.9,0.9,0.9}{
            \makecell[{{p{\linewidth}}}]{
                \texttt{\tiny{[GM$\langle$P2]}}
                \texttt{VRAAG: Zijn er rijen met een ander aantal lege vakken (▢) aan het einde van die rijen, naast de reeds genoemde rijen met drie X{-}tekens?} \\
            }
        }
    }
     \\ \\

    \theutterance \stepcounter{utterance}  
    & & & \multicolumn{2}{p{0.3\linewidth}}{
        \cellcolor[rgb]{0.9,0.9,0.9}{
            \makecell[{{p{\linewidth}}}]{
                \texttt{\tiny{[GM$|$GM]}}
                \texttt{continue} \\
            }
        }
    }
    & & \\ \\

    \theutterance \stepcounter{utterance}  
    & & \multicolumn{4}{p{0.6\linewidth}}{
        \cellcolor[rgb]{0.9,0.9,0.9}{
            \makecell[{{p{\linewidth}}}]{
                \texttt{\tiny{[P1$\langle$GM]}}
                \texttt{ANTWOORD: Ik kan geen kolommen met X{-}tekens visualiseren, aangezien ik alleen de rasterbeschrijving heb. Ik kan dus geen antwoord geven op vragen over kolommen zonder verdere informatie.} \\
\texttt{VRAAG: Zijn er rijen met een ander aantal lege vakken (▢) aan het einde van die rijen, naast de reeds genoemde rijen met drie X{-}tekens?Begin je antwoord met "ANTWOORD:" en voeg niets anders toe.} \\
            }
        }
    }
    & \\ \\

    \theutterance \stepcounter{utterance}  
    & \multicolumn{4}{p{0.6\linewidth}}{
        \cellcolor[rgb]{0.9,0.9,0.9}{
            \makecell[{{p{\linewidth}}}]{
                \texttt{\tiny{[P1$\rangle$GM]}}
                \texttt{ANTWOORD: Ik kan geen antwoord geven zonder eerst te weten hoe jouw raster precies is gestructureerd na de reeds genoemde rijen. Geef meer details over het aantal lege vakken (▢) in die andere rijen, zodat ik een juist antwoord kan geven.} \\
            }
        }
    }
    & & \\ \\

    \theutterance \stepcounter{utterance}  
    & & & \multicolumn{2}{p{0.3\linewidth}}{
        \cellcolor[rgb]{0.9,0.9,0.9}{
            \makecell[{{p{\linewidth}}}]{
                \texttt{\tiny{[GM$|$GM]}}
                \texttt{continue} \\
            }
        }
    }
    & & \\ \\

    \theutterance \stepcounter{utterance}  
    & & \multicolumn{4}{p{0.6\linewidth}}{
        \cellcolor[rgb]{0.9,0.9,0.9}{
            \makecell[{{p{\linewidth}}}]{
                \texttt{\tiny{[P1$\langle$GM]}}
                \texttt{Stel nu een vraag om nieuwe aspecten van mijn raster te ontdekken die anders kunnen zijn dan jouw raster. Begin met "VRAAG:" en voeg niets anders toe.} \\
            }
        }
    }
    & \\ \\

    \theutterance \stepcounter{utterance}  
    & \multicolumn{4}{p{0.6\linewidth}}{
        \cellcolor[rgb]{0.9,0.9,0.9}{
            \makecell[{{p{\linewidth}}}]{
                \texttt{\tiny{[P1$\rangle$GM]}}
                \texttt{VRAAG: Zijn er rijen waarin de lege vakken (▢) niet aan het begin van de rij staan?} \\
            }
        }
    }
    & & \\ \\

    \theutterance \stepcounter{utterance}  
    & & & \multicolumn{2}{p{0.3\linewidth}}{
        \cellcolor[rgb]{0.9,0.9,0.9}{
            \makecell[{{p{\linewidth}}}]{
                \texttt{\tiny{[GM$|$GM]}}
                \texttt{continue} \\
            }
        }
    }
    & & \\ \\

    \theutterance \stepcounter{utterance}  
    & & \multicolumn{4}{p{0.6\linewidth}}{
        \cellcolor[rgb]{0.9,0.9,0.9}{
            \makecell[{{p{\linewidth}}}]{
                \texttt{\tiny{[GM$\rangle$P2]}}
                \texttt{ANTWOORD: Ik kan geen antwoord geven zonder eerst te weten hoe jouw raster precies is gestructureerd na de reeds genoemde rijen. Geef meer details over het aantal lege vakken (▢) in die andere rijen, zodat ik een juist antwoord kan geven.} \\
\texttt{VRAAG: Zijn er rijen waarin de lege vakken (▢) niet aan het begin van de rij staan?Begin je antwoord met "ANTWOORD:" en voeg niets anders toe.} \\
            }
        }
    }
    & \\ \\

    \theutterance \stepcounter{utterance}  
    & & & \multicolumn{4}{p{0.6\linewidth}}{
        \cellcolor[rgb]{0.9,0.9,0.9}{
            \makecell[{{p{\linewidth}}}]{
                \texttt{\tiny{[GM$\langle$P2]}}
                \texttt{ANTWOORD: Ik kan geen antwoord geven zonder meer context over de positie van de lege vakken in jouw raster.} \\
            }
        }
    }
     \\ \\

    \theutterance \stepcounter{utterance}  
    & & & \multicolumn{2}{p{0.3\linewidth}}{
        \cellcolor[rgb]{0.9,0.9,0.9}{
            \makecell[{{p{\linewidth}}}]{
                \texttt{\tiny{[GM$|$GM]}}
                \texttt{continue} \\
            }
        }
    }
    & & \\ \\

    \theutterance \stepcounter{utterance}  
    & & \multicolumn{4}{p{0.6\linewidth}}{
        \cellcolor[rgb]{0.9,0.9,0.9}{
            \makecell[{{p{\linewidth}}}]{
                \texttt{\tiny{[GM$\rangle$P2]}}
                \texttt{Stel nu een vraag om nieuwe aspecten van mijn raster te ontdekken die anders kunnen zijn dan jouw raster. Begin met "VRAAG:" en voeg niets anders toe.} \\
            }
        }
    }
    & \\ \\

    \theutterance \stepcounter{utterance}  
    & & & \multicolumn{4}{p{0.6\linewidth}}{
        \cellcolor[rgb]{0.9,0.9,0.9}{
            \makecell[{{p{\linewidth}}}]{
                \texttt{\tiny{[GM$\langle$P2]}}
                \texttt{VRAAG: Zijn er rijen met een andere volgorde van lege vakken (▢) en X{-}tekens, zoals bijvoorbeeld een reeks met een X in een positie die niet overeenkomt met jouw raster?} \\
            }
        }
    }
     \\ \\

    \theutterance \stepcounter{utterance}  
    & & & \multicolumn{2}{p{0.3\linewidth}}{
        \cellcolor[rgb]{0.9,0.9,0.9}{
            \makecell[{{p{\linewidth}}}]{
                \texttt{\tiny{[GM$|$GM]}}
                \texttt{continue} \\
            }
        }
    }
    & & \\ \\

    \theutterance \stepcounter{utterance}  
    & & \multicolumn{4}{p{0.6\linewidth}}{
        \cellcolor[rgb]{0.9,0.9,0.9}{
            \makecell[{{p{\linewidth}}}]{
                \texttt{\tiny{[P1$\langle$GM]}}
                \texttt{ANTWOORD: Ik kan geen antwoord geven zonder meer context over de positie van de lege vakken in jouw raster.} \\
\texttt{VRAAG: Zijn er rijen met een andere volgorde van lege vakken (▢) en X{-}tekens, zoals bijvoorbeeld een reeks met een X in een positie die niet overeenkomt met jouw raster?Begin je antwoord met "ANTWOORD:" en voeg niets anders toe.} \\
            }
        }
    }
    & \\ \\

    \theutterance \stepcounter{utterance}  
    & \multicolumn{4}{p{0.6\linewidth}}{
        \cellcolor[rgb]{0.9,0.9,0.9}{
            \makecell[{{p{\linewidth}}}]{
                \texttt{\tiny{[P1$\rangle$GM]}}
                \texttt{ANTWOORD: Ik kan geen antwoord geven zonder te weten hoe jouw raster precies eruitziet ten opzichte van mijn rasterbeschrijving. Geef een voorbeeld van de volgorde van lege vakken (▢) en X{-}tekens in jouw raster om een juist antwoord te kunnen geven.} \\
            }
        }
    }
    & & \\ \\

    \theutterance \stepcounter{utterance}  
    & & & \multicolumn{2}{p{0.3\linewidth}}{
        \cellcolor[rgb]{0.9,0.9,0.9}{
            \makecell[{{p{\linewidth}}}]{
                \texttt{\tiny{[GM$|$GM]}}
                \texttt{continue} \\
            }
        }
    }
    & & \\ \\

    \theutterance \stepcounter{utterance}  
    & & \multicolumn{4}{p{0.6\linewidth}}{
        \cellcolor[rgb]{0.9,0.9,0.9}{
            \makecell[{{p{\linewidth}}}]{
                \texttt{\tiny{[P1$\langle$GM]}}
                \texttt{Stel nu een vraag om nieuwe aspecten van mijn raster te ontdekken die anders kunnen zijn dan jouw raster. Begin met "VRAAG:" en voeg niets anders toe.} \\
            }
        }
    }
    & \\ \\

    \theutterance \stepcounter{utterance}  
    & \multicolumn{4}{p{0.6\linewidth}}{
        \cellcolor[rgb]{0.9,0.9,0.9}{
            \makecell[{{p{\linewidth}}}]{
                \texttt{\tiny{[P1$\rangle$GM]}}
                \texttt{VRAAG: Zijn er rijen met een onregelmatig aantal lege vakken (▢) of X{-}tekens die niet voorkomen in mijn rasterbeschrijving?} \\
            }
        }
    }
    & & \\ \\

    \theutterance \stepcounter{utterance}  
    & & & \multicolumn{2}{p{0.3\linewidth}}{
        \cellcolor[rgb]{0.9,0.9,0.9}{
            \makecell[{{p{\linewidth}}}]{
                \texttt{\tiny{[GM$|$GM]}}
                \texttt{continue} \\
            }
        }
    }
    & & \\ \\

    \theutterance \stepcounter{utterance}  
    & & \multicolumn{4}{p{0.6\linewidth}}{
        \cellcolor[rgb]{0.9,0.9,0.9}{
            \makecell[{{p{\linewidth}}}]{
                \texttt{\tiny{[GM$\rangle$P2]}}
                \texttt{ANTWOORD: Ik kan geen antwoord geven zonder te weten hoe jouw raster precies eruitziet ten opzichte van mijn rasterbeschrijving. Geef een voorbeeld van de volgorde van lege vakken (▢) en X{-}tekens in jouw raster om een juist antwoord te kunnen geven.} \\
\texttt{VRAAG: Zijn er rijen met een onregelmatig aantal lege vakken (▢) of X{-}tekens die niet voorkomen in mijn rasterbeschrijving?Begin je antwoord met "ANTWOORD:" en voeg niets anders toe.} \\
            }
        }
    }
    & \\ \\

    \theutterance \stepcounter{utterance}  
    & & & \multicolumn{4}{p{0.6\linewidth}}{
        \cellcolor[rgb]{0.9,0.9,0.9}{
            \makecell[{{p{\linewidth}}}]{
                \texttt{\tiny{[GM$\langle$P2]}}
                \texttt{ANTWOORD: Ik kan geen antwoord geven zonder informatie over het aantal lege vakken (▢) en X{-}tekens in jouw raster.} \\
            }
        }
    }
     \\ \\

    \theutterance \stepcounter{utterance}  
    & & & \multicolumn{2}{p{0.3\linewidth}}{
        \cellcolor[rgb]{0.9,0.9,0.9}{
            \makecell[{{p{\linewidth}}}]{
                \texttt{\tiny{[GM$|$GM]}}
                \texttt{continue} \\
            }
        }
    }
    & & \\ \\

    \theutterance \stepcounter{utterance}  
    & & \multicolumn{4}{p{0.6\linewidth}}{
        \cellcolor[rgb]{0.9,0.9,0.9}{
            \makecell[{{p{\linewidth}}}]{
                \texttt{\tiny{[GM$\rangle$P2]}}
                \texttt{Neem nu een beslissing. Wat denk je: zijn jouw raster en het andere beschreven raster hetzelfde? Schrijf "BESLISSING: hetzelfde raster" als je denkt dat het hetzelfde raster is of "BESLISSING: verschillende rasters" als je denkt dat het verschillende rasters zijn. Voeg niets anders toe.} \\
            }
        }
    }
    & \\ \\

    \theutterance \stepcounter{utterance}  
    & & & \multicolumn{4}{p{0.6\linewidth}}{
        \cellcolor[rgb]{0.9,0.9,0.9}{
            \makecell[{{p{\linewidth}}}]{
                \texttt{\tiny{[GM$\langle$P2]}}
                \texttt{BESLISSING: verschillende rasters} \\
            }
        }
    }
     \\ \\

    \theutterance \stepcounter{utterance}  
    & & & \multicolumn{2}{p{0.3\linewidth}}{
        \cellcolor[rgb]{0.9,0.9,0.9}{
            \makecell[{{p{\linewidth}}}]{
                \texttt{\tiny{[GM$|$GM]}}
                \texttt{continue} \\
            }
        }
    }
    & & \\ \\

    \theutterance \stepcounter{utterance}  
    & & & \multicolumn{2}{p{0.3\linewidth}}{
        \cellcolor[rgb]{0.9,0.9,0.9}{
            \makecell[{{p{\linewidth}}}]{
                \texttt{\tiny{[GM$|$GM]}}
                \texttt{success} \\
            }
        }
    }
    & & \\ \\

    \theutterance \stepcounter{utterance}  
    & & \multicolumn{4}{p{0.6\linewidth}}{
        \cellcolor[rgb]{0.9,0.9,0.9}{
            \makecell[{{p{\linewidth}}}]{
                \texttt{\tiny{[P1$\langle$GM]}}
                \texttt{ANTWOORD: Ik kan geen antwoord geven zonder informatie over het aantal lege vakken (▢) en X{-}tekens in jouw raster.} \\
\texttt{Neem nu een beslissing. Wat denk je: zijn jouw raster en het andere beschreven raster hetzelfde? Schrijf "BESLISSING: hetzelfde raster" als je denkt dat het hetzelfde raster is of "BESLISSING: verschillende rasters" als je denkt dat het verschillende rasters zijn. Voeg niets anders toe.} \\
            }
        }
    }
    & \\ \\

    \theutterance \stepcounter{utterance}  
    & \multicolumn{4}{p{0.6\linewidth}}{
        \cellcolor[rgb]{0.9,0.9,0.9}{
            \makecell[{{p{\linewidth}}}]{
                \texttt{\tiny{[P1$\rangle$GM]}}
                \texttt{BESLISSING: verschillende rasters} \\
            }
        }
    }
    & & \\ \\

    \theutterance \stepcounter{utterance}  
    & & & \multicolumn{2}{p{0.3\linewidth}}{
        \cellcolor[rgb]{0.9,0.9,0.9}{
            \makecell[{{p{\linewidth}}}]{
                \texttt{\tiny{[GM$|$GM]}}
                \texttt{continue} \\
            }
        }
    }
    & & \\ \\

    \theutterance \stepcounter{utterance}  
    & & & \multicolumn{2}{p{0.3\linewidth}}{
        \cellcolor[rgb]{0.9,0.9,0.9}{
            \makecell[{{p{\linewidth}}}]{
                \texttt{\tiny{[GM$|$GM]}}
                \texttt{success} \\
            }
        }
    }
    & & \\ \\

\end{supertabular}
}

\end{document}
